\chapter{[\textit{Title of Your Main Contribution}]}\label{chap:main-chapter}

This chapter contains the main theoretical (and possibly technical) contribution
of the thesis.  It should present a coherent body of new results, organized
into sections.

For a theory-focused thesis: introduce new notions, state lemmas/theorems,
provide proofs, and illustrate results with examples.
For an implementation-focused thesis: additionally describe the algorithm,
key design decisions, and data structures.

The opening paragraph of this chapter should cover the following points:
\begin{enumerate}
    \item State what new objects, criteria, or techniques this chapter introduces.
    \item Introduce a running example that will be used throughout to illustrate
          each new definition and result.
    \item Outline the chapter structure by referencing the main sections.
\end{enumerate}

%-----------------------------------------------------------------------
\section{[\textit{Central Definitions / The New Framework}]}\label{sec:main-definitions}
%-----------------------------------------------------------------------

% Introduce the key new notions.
% Follow the style of \cref{chap:preliminaries}: state each definition formally,
% add a brief intuitive explanation, and accompany it with at least one example.
% Move from simple to complex.

For self-introduced notions it is even more important to explain everything
and used multiple examples to illustrate everything.
Use one running example that can be used through the whole section 
and maybe some additional small examples to highlight special cases.

\begin{Definition}{[\textit{Name of Your First Definition}]}{key1}

\end{Definition}

\begin{example}\label{ex:key1}

\end{example}

% Add further Definition boxes as needed, each followed by an explanation and example.

%-----------------------------------------------------------------------
\section{[\textit{Main Results}]}\label{sec:main-results}
%-----------------------------------------------------------------------

Each result in this section should be presented as follows:
\begin{enumerate}
    \item State and prove all auxiliary lemmas before the main theorem.
    \item State the main theorem.  Typical statements are:
          \begin{itemize}
              \item \textbf{Soundness:} If the technique returns \textsc{Yes},
                    then the property holds.
              \item \textbf{Completeness:} If the property holds, the technique
                    returns \textsc{Yes}.
              \item \textbf{Equivalence:} The new criterion holds if and only if
                    an existing criterion holds.
          \end{itemize}
    \item Provide a complete proof, structured into numbered cases for long
          arguments.  Reference definitions and earlier results by label.
    \item Apply the theorem to the running example from \cref{ex:key1},
          walking through every step explicitly.
\end{enumerate}

\begin{Lemma}{[\textit{Name of Auxiliary Lemma}]}{auxiliary}

\end{Lemma}

\begin{proof}

\end{proof}

\begin{Theorem}{[\textit{Name of Main Theorem, e.g., Soundness}]}{main}

\end{Theorem}

\begin{proof}

\end{proof}

\begin{example}\label{ex:main}
    Apply \Cref{thm:main} to the running example from \cref{ex:key1}.
    Walk through every step of the argument explicitly.
\end{example}